\section{Introduction}
\label{sec:introduction}

Science has always had an objective function. It has never been written
down in a single place. This paper writes it down.

The objective function is:
\[
  \text{minimize} \quad L(H) + L(D \mid H)
\]
where $H$ is a hypothesis, $D$ is observed data, $L(H)$ is the
description length of the hypothesis, and $L(D \mid H)$ is the
description length of the data given the hypothesis. This is the
Minimum Description Length (MDL) principle---the computable
approximation to Solomonoff induction, which is Bayesian inference
with the Kolmogorov/Solomonoff prior. It is not a new proposal. It is
a formal statement of what science has always been doing informally.

This paper makes four contributions.

\textbf{A synthesis of the history of philosophy of science.} The
history of philosophy of science from Hume through Popper, Kuhn,
Lakatos, Feyerabend, and the Bayesian tradition is a single long
argument about one problem: Hume's problem of induction. Every major
movement is a partial approximation to the framework above. We trace
this arc and show that each movement captured one component of the
full solution without being able to state it. Section~\ref{sec:history}
develops this argument. Section~\ref{sec:connections} identifies ten
cases where ideas from independent traditions---philosophy,
computability theory, information theory, neuroscience, machine
learning---turn out to be the same idea expressed in different language.

\textbf{The QBBN as a universal coding scheme for logical hypotheses.}
MDL requires a coding scheme: a way of assigning description lengths to
hypotheses. The Kolmogorov prior gives the theoretically optimal
scheme, but it is uncomputable. Every practical coding scheme covers
only a restricted hypothesis class. The Quantified Boolean Bayesian
Network (QBBN)~\cite{coppola2024qbbn} is the best known candidate for
a universal practical coding scheme: it is complete over first-order
logic (covering all of mathematics and science), generative (assigning
likelihoods to arbitrary proposition datasets), and tractable (belief
propagation runs in polynomial time). Moreover, the
\texttt{shannon}~paper~\cite{coppola2026} proves that transformers
implement exact belief propagation on QBBN-structured factor graphs,
connecting the coding scheme to the architecture that currently
dominates machine learning. Section~\ref{sec:coding} develops this
argument.

\textbf{Two empirical demonstrations.} The framework is demonstrated
empirically in two domains. In the domain of code and programs
(Section~\ref{sec:experiments-code}), we demonstrate compression
ratios, MDL model selection, the replication crisis as a missing
$L(H)$ term, the search/verify asymmetry, and the Karpathy
gap---current AutoResearch systems minimize $L(D \mid H)$ without
$L(H)$, and the missing term accounts for a 65.2\% improvement on
benchmark problems. In the domain of logical propositions
(Section~\ref{sec:experiments-logic}), we use the QBBN coding scheme
to show that conjunctive structure wins over disjunctive, that
theories are amortized over data (8 bits of theory saves 510 bits at
$n = 1000$), and that MDL selects correct theory complexity as a
function of data regularity, with a discrete crossover that
instantiates Kuhn's paradigm shift mechanism at the bit level.

\textbf{Theorem proving is science; the objective function is AGI.}
The critics who say ``automated theorem proving is not AGI because it
is not science'' are relying on an implicit theory of science that has
no formal content. Under the correct theory, the verification oracle
is interchangeable: a particle accelerator, a type-checker, and an
MDL score calculation are all instances of the same abstract
structure. A system with the right objective function is doing science
regardless of substrate. Section~\ref{sec:agi} develops this argument
and proposes a formal definition of AGI grounded in the framework.

\subsection*{Related Work}

The Bayesian/Kolmogorov framework synthesized here draws on
Solomonoff's universal prior~\cite{solomonoff1964}, Kolmogorov
complexity~\cite{kolmogorov1965}, Rissanen's MDL
principle~\cite{rissanen1978}, and the Bayesian tradition in
philosophy of science~\cite{howson2006}. The connection between
paradigm shifts and MDL phase transitions extends
Kuhn~\cite{kuhn1962} and Lakatos~\cite{lakatos1978}. The
Bayesian brain and predictive coding framework draws on Rao and
Ballard~\cite{rao1999} and Friston~\cite{friston2005}. The QBBN
coding scheme is developed in~\cite{coppola2024qbbn}. The proof that
transformers implement exact belief propagation on QBBN-structured
factor graphs is in~\cite{coppola2026}. The AutoResearch system
discussed in the Karpathy gap analysis is~\cite{karpathy2026}.

The \texttt{auto-research} and \texttt{mdl-science} repositories
contain all code and data for the empirical demonstrations. All
results are reproducible by following the instructions in those
repositories.